%------------------------------------------------------------------------------
% 4.6 تصميم واجهات برمجة التطبيقات (APIs)
% المالك: BE (مهمة جانت ٢.٢.٥ — 06/11 إلى 09/11/2026)
% الحجم المستهدف: 2-3 صفحات
%
% عناصر مطلوبة:
%   [ ] 4.6.1 الأسلوب المعتمد: REST فوق HTTP/JSON + بث الأحداث من الخادم
%       (SSE) لمسار المحادثة لعرض الإجابة تدريجياً — وهو الجوهر التقني
%       لتجربة الاستجابة السريعة الموعودة في المتطلبات غير الوظيفية
%   [ ] 4.6.2 كتالوج نقاط النهاية (tab:api-endpoints): الطريقة / المسار /
%       الغرض / الدور المسموح — نحو 8-12 نقطة تغطي: المصادقة، المحادثة
%       (والبث)، التقارير والتصدير، إدارة المستخدمين والأدوار، العروض
%       الذكية، سجل التدقيق، التهيئة الذكية، فحص الصحة
%   [ ] 4.6.3 معالجة الأخطاء: صيغة الاستجابة الموحدة (رمز خطأ داخلي ثابت،
%       رسالة عربية للمستخدم، تفاصيل تشخيصية للمطورين) + رموز حالات HTTP
%       المعتمدة لكل عائلة أخطاء (صلاحيات / تحقق من المدخلات / حدود المعدل /
%       أخطاء الخادم)
%
% مرجع الدقة: مسارات API الفعلية في مستودع Bayan (backend/src/api) — ابنِ
%   الكتالوج من الكود الموجود لا من التصور
%------------------------------------------------------------------------------
\section{تصميم واجهات برمجة التطبيقات (\EN{APIs})}
\label{sec:api-design}
~

\subsection{الأسلوب المعتمد}
\label{subsec:api-style}
~

\subsection{كتالوج نقاط النهاية}
\label{subsec:api-catalog}
~

\subsection{معالجة الأخطاء وصيغ الاستجابة الموحدة}
\label{subsec:api-errors}
~
